\chapter{Configurations}
\label{ch:configurations}

A configuration is a named recipe --- one complete set of the 31 bonding
parameters. You will normally keep one per job: per wire type, per device, per
pad metallisation.

The machine stores up to 32 configurations in total.

\section{Configurations belong to a mode}

Each configuration belongs to one bonding mode, and the list only shows the
configurations for the mode you are viewing. Use $\leftarrow$ and $\rightarrow$
on the configuration list to switch modes.

\begin{wbnote}
This means a Semi Auto recipe and a Manual recipe can share a name without
clashing --- they are different configurations in different lists.
\end{wbnote}

\section{DEFAULT}

Every mode has a configuration called \menuitem{DEFAULT}, created automatically
the first time the machine runs. It is your fallback.

\begin{wbnote}
\menuitem{DEFAULT} cannot be deleted --- the machine answers
\msg{Cannot delete DEFAULT} --- and pressing \keycap{SAVE} while it is loaded
opens Save As instead of overwriting it. This is deliberate: it guarantees
there is always one known-good starting point to fall back to.
\end{wbnote}

To use \menuitem{DEFAULT} as a starting point, load it, make your changes, then
press \keycap{SAVE} and give the result a new name.

\section{Creating a configuration}

\begin{enumerate}[leftmargin=*]
\item From the configuration list, select the mode you want with $\leftarrow$
  and $\rightarrow$.
\item Press \keycap{ADD}.
\item Type a name and press \keycap{ENTER}.
\end{enumerate}

The new configuration is created with factory default values
(Appendix~\ref{app:defaults}) and loaded ready for editing.

\section{Loading}

Move the \texttt{*} to the configuration you want and press \keycap{ENTER} or
\keycap{LOAD}. The parameter editor opens with it loaded.

\section{Saving}

Press \keycap{SAVE} on the Parameter Edit page. The machine confirms with
\msg{Configuration saved}.

\begin{wbwarning}
Nothing is saved automatically. Edits live only in memory until you press
\keycap{SAVE}, and are lost at power-down. Loading another configuration
discards them too, without warning.
\end{wbwarning}

If the loaded configuration is \menuitem{DEFAULT}, \keycap{SAVE} opens Save As
so you can store the result under a new name.

\section{Deleting}

From the configuration list, move the \texttt{*} to the configuration and press
\keycap{ESC/DEL}. Answer \menuitem{Y} and press \keycap{ENTER}.

\begin{wbwarning}
There is no undo. A deleted configuration is gone.
\end{wbwarning}

\section{When storage is full}

With 32 configurations stored, creating another fails with \msg{Add failed} and
saving a new name fails with \msg{Save failed}. Delete something you no longer
need.

\TODO{If a house convention for naming configurations exists --- by wire
diameter, device type, customer or job number --- describe it here, with
worked examples. A consistent scheme matters because names are limited to ten
characters.}

\section{Recommended practice}

\TODO{Write the site's recipe-control practice: who is permitted to change a
qualified configuration, how changes are recorded, how a recipe is re-qualified
after a tool or wire change, and how configurations are backed up given that
the machine has no export function.}
